\chapter{Conclusion}


\section{Thesis Summary}
This thesis investigated automated prediction of patient-level mucus plug burden from lung CT scans. The task was formulated as weakly supervised regression, each CT examination had one patient-level mucus-burden score, while the model operated on selected axial slices without slice-level mucus annotations. This made the problem different from direct plug detection or segmentation. The goal was not to build a clinically ready system, but to test whether a constrained slice-based pipeline could learn useful patient-level signals from the available data.

The implemented pipeline combined stored CT volume loading, intensity handling, representative slice selection, CT input representation, CNN feature extraction, optional ORB-BoVW feature fusion, slice-level prediction, and patient-level aggregation. The main internal evaluation used five-fold cross-validation on the internal mucus-score cohort, while the transfer experiments used MosMed as a separate CT dataset for studying representation transfer and domain shift. Because MosMed has a COVID-related severity target rather than a mucus-burden target, it was not treated as external validation of mucus prediction.

The results answer the research questions in a cautious but useful way. For RQ1, the internal experiments showed that patient-level prediction is possible to approach with the proposed weakly supervised pipeline, but the low and sometimes negative \(R^2\) values also show that the task remains difficult. For RQ2, representative slice selection performed better than uniform slice sampling in the baseline comparison, suggesting that the way slices are selected matters when the full CT volume is not processed end-to-end. For RQ3, the three-channel high-pass input and ORB-BoVW fusion were associated with the strongest internal results, with the ResNet18 three-channel fusion model performing best among the reported internal configurations. For RQ4, DenseNet121 did not consistently outperform ResNet18, even though it gave slightly better error metrics in one three-channel setting. For RQ5, the transfer experiments showed that fine-tuning was more useful than zero-shot transfer, while also showing the limits of direct transfer when the target definition and domain change.

These findings support the main idea of the thesis, that predicting mucus burden from CT scans at the patient level can be explored using a slice-based weakly supervised pipeline, and that some design choices seem helpful in this experimental setup. At the same time, the results should still be seen as exploratory. The cohort is small, the labels are only available at the patient level, and the transfer experiments are not a substitute for validation on an independent dataset with mucus scores.


\section{Main Contributions}
The first contribution is the formulation of mucus plug burden prediction as a patient-level weakly supervised regression problem. This framing is important because it keeps the prediction target aligned with the available clinical score while acknowledging that the image evidence is distributed across a volumetric CT examination. It also makes clear that slice-level outputs from the model are intermediate signals, not verified mucus-localization results.

The second contribution is the implementation of a complete and reproducible slice-based CT pipeline. The pipeline separates the main stages of the problem: preprocessing, representative slice selection strategy, input representation, feature extraction, aggregation, and model selection. This separation made it possible to test the role of different parts instead of treating the model as one whole system.

The third contribution of this thesis is comparing different design choices when only a small amount of data is available. For example, representative slice selection was compared with uniform sampling, two-window input was compared with a three-channel high-pass method, and ResNet18 was compared with DenseNet121. The study also tested ORB-BoVW fusion as a handcrafted feature branch. These comparisons are not perfect ablation studies, but they still help show which choices gave better results in this thesis.

The fourth contribution is the transfer-learning analysis using MosMed. These experiments did not validate mucus prediction externally, but they did give useful information about source--target adaptation between CT tasks. The difference between zero-shot transfer and fine-tuning showed that source-trained CT representations can be helpful, but that direct reuse across a different target definition is limited.


\section{Future Work}
The most important next step is evaluation on a larger mucus-scored cohort. With only 32 internal cases, each validation fold is small, and individual patients can strongly affect the reported metrics. A larger cohort would make the results more stable and would make it easier to test whether the observed patterns for slice selection, high-pass input, feature fusion, and backbone choice remain consistent.

External validation should also be performed on an independent dataset with a comparable mucus-burden target. The MosMed experiments were useful for studying transfer and domain shift, but they cannot answer whether the model generalizes to a separate mucus-scored population. A true external validation dataset would be needed before making stronger claims about generalization.

Future work should also include more detailed annotations if they become available. Slice-level or region-level mucus labels would make it possible to evaluate whether high-scoring slices correspond to actual mucus findings. This would strengthen interpretability and would help separate patient-level prediction from localization, which this thesis could not do.

Several methodological directions are also worth extending. A more controlled aggregation study could compare mean, top-\(k\), hybrid, and attention-based pooling in a cleaner way. The input-representation variants could be evaluated more systematically, including different high-pass settings and other frequency-based channels. The ORB-BoVW branch should also be tested in larger data to see whether handcrafted local features remain useful when the CNN has more examples to learn from. Finally, models that use more spatial context, such as 2.5D inputs, short slice sequences, or volumetric models, may be useful if enough data are available to support them.


\section{Concluding Remarks}
This thesis shows that automated patient-level mucus-burden prediction from CT is a meaningful but difficult machine-learning problem. The best internal results came from a pipeline that combined representative slice selection strategy, a three-channel high-pass CT representation, ResNet18 feature extraction, ORB-BoVW fusion, and patient-level aggregation. These findings suggest that careful input construction and complementary feature representations matter in a small weakly supervised CT setting.

The work also shows why the problem should be treated cautiously. Low and negative \(R^2\) values in several settings, the small internal cohort, patient-level-only labels, and the absence of external mucus-specific validation all limit the strength of the conclusions. The transfer experiments further show that CT representations do not transfer reliably without adaptation when the dataset and target definition change.

The final takeaway is therefore balanced. The proposed pipeline is not a clinical mucus-plug detection system, and it does not prove that the model localizes mucus plugs. What it does provide is a reproducible experimental framework and evidence that patient-level mucus-burden prediction from CT can benefit from thoughtful slice selection, input representation, feature fusion, and fine-tuning. This makes it a useful starting point for larger, better-annotated, and externally validated studies.

